% Auto-generated by scripts/run_ew_portfolio.py -- do not edit by hand.
% Requires: \usepackage{booktabs}
\begin{table}[t]
\centering
\small
\begin{tabular}{lrrrr}
\toprule
Agent (EW, \%) & bull\_2019 & crisis\_2020 & recovery\_2021 & hi\_rates\_2022 \\
\midrule
Buy \& Hold & 17.2 & -22.4 & 10.0 & 14.6 \\
MACD(12,26,9) & 4.9 & 6.9 & 8.8 & 14.6 \\
SMA(50/200) & 11.2 & -26.0 & 8.0 & -11.0 \\
Momentum(60) & 4.5 & -4.5 & 3.1 & 12.2 \\
Random ($N{=}100$ mean) & 9.7 & -12.3 & 6.2 & 6.3 \\
\bottomrule
\end{tabular}
\caption{\textbf{Equal-weight portfolio returns by rule-based agent and regime.} Each entry is the naive equal-weight (1/5) portfolio total return across the five tickers (PETR4, VALE3, ITUB4, BBDC4, \texttt{\^{}BVSP}), i.e.\ the mean of the per-ticker cumulative returns, over the four dated regime windows. All strategies are long-only, binary (all-in/all-cash), and frictionless; the trend strategies (MACD, SMA, Momentum) use a pre-window warmup buffer so their indicators are warm at each window's first bar (no look-ahead). The Random row is the $N{=}100$-seed mean per ticker from Table~\ref{tab:random-n100}, i.e.\ the expected uninformed return rather than a single noisy seed. \emph{LLM excluded: sequential multi-seed evaluation prevents proper simultaneous portfolio construction.} These are classical baselines only; the multi-agent LLM pipeline is evaluated separately.}
\label{tab:ew-baselines}
\end{table}
